\subsection{Zyklische Ablaufsteuerung}
Nach der Initialisierung wird das System in einer Endlosschleife betrieben. 
Innerhalb dieser Schleife werden zyklisch Webanfragen verarbeitet, Zeitinformationen aktualisiert sowie Anzeige- und Sensorfunktionen ausgeführt. 
Die zeitliche Steuerung erfolgt in festen Intervallen, wodurch eine gleichmäßige und deterministische Aktualisierung gewährleistet wird. 
Der vollständige Ablauf ist in Abb.~\ref{fig:hauptprogramm} dargestellt.

\begin{figure}[H]
\centering
\begin{tikzpicture}[
    node distance=0.45cm,
    startstop/.style={rectangle, rounded corners, minimum width=4cm, minimum height=0.8cm, text centered, draw=black},
    process/.style={rectangle, minimum width=4.6cm, minimum height=0.8cm, text centered, draw=black},
    decision/.style={diamond, aspect=2.4, minimum width=4cm, minimum height=1cm, text centered, draw=black},
    io/.style={trapezium, trapezium left angle=70, trapezium right angle=110, minimum width=4.6cm, minimum height=0.8cm, text centered, draw=black},
    arrow/.style={thick, ->, >=stealth}
]

\node (start) [startstop] {Start};
\node (setup) [process, below=of start] {setup ausführen};
\node (loop) [process, below=of setup] {loop dauerhaft wiederholen};
\node (web) [process, below=of loop] {Webserver-Anfragen bearbeiten};
\node (check500) [decision, below=of web] {500 ms vergangen?};
\node (time) [process, below=of check500] {Zeit abfragen};
\node (valid) [decision, below=of time] {Zeit gültig?};
\node (daily) [process, below=of valid] {Zeit aktualisieren};
\node (weather) [process, below=of daily] {Wetteraktualisierung abrufen};
\node (sensor) [process, below=of weather] {Innenraumsensordaten abrufen};
\node (wordclock) [io, below=of sensor] {Hauptanzeige anzeigen};
\node (second) [io, below=of wordclock] {Sekunden anzeigen};
\node (third) [io, below=of second] {Zusatzfunktionen anzeigen};

\draw [arrow] (start) -- (setup);
\draw [arrow] (setup) -- (loop);
\draw [arrow] (loop) -- (web);
\draw [arrow] (web) -- (check500);

\draw [arrow] (check500) -- node[right] {Ja} (time);
\draw [arrow] (time) -- (valid);
\draw [arrow] (valid) -- node[right] {Ja} (daily);
\draw [arrow] (daily) -- (weather);
\draw [arrow] (weather) -- (sensor);
\draw [arrow] (sensor) -- (wordclock);
\draw [arrow] (wordclock) -- (second);
\draw [arrow] (second) -- (third);

\draw [arrow] (check500.east) -- ++(1.8,0)
    node[midway, above] {Nein}
    |- (loop.east);

\draw [arrow] (valid.east) -- ++(3.0,0)
    node[midway, above] {Nein}
    |- (loop.east);

\draw [arrow] (third.east) -- ++(3.0,0)
    |- (loop.east);

\end{tikzpicture}
\caption[Schematischer Programmablaufplan des zyklischen Hauptprogramms]{Schematischer Programmablaufplan des sich zyklisch wiederholenden Hauptprogramms der WordClock.}
\label{fig:hauptprogramm}
\end{figure}